\section*{Extensions and Limitations}

One limitation is that headquarters county does not always show where a firm really operates. A firm can be headquartered outside the FEMA area but have factories or shops inside it, or the other way round. In that case I put firms in the wrong group. This kind of mistake normally makes the effect look smaller than it really is, so my $b$ would go toward zero. A second problem is that cutting R\&D automatically leaves more cash in the firm. So the two results alone cannot fully prove that managers acted out of precaution.

As an extension I will test if the effect is stronger for firms that had problems financing before Helene. I use a triple interaction $\mathrm{Treated}_i \cdot \mathrm{Post}_t \cdot \mathrm{Constrained}_i$, where $\mathrm{Constrained}_i$ is 1 for firms with high financing constraints in 2023 and 0 if not. I measure constraints with firm size and leverage before the shock. If constrained firms cut R\&D more and build up more cash than others, this supports the self-insurance idea, because these firms cannot easily borrow from banks after a shock.
